\documentclass[10pt]{article}
\usepackage[T1]{fontenc}
\usepackage[utf8]{inputenc}
\usepackage[french]{babel}
\usepackage{amsmath,amssymb,mathrsfs,mathpazo}
\usepackage{geometry}
\geometry{a4paper,margin=2cm}
\usepackage{graphicx,array,tabularx,tabularray,colortbl,multirow,multicol,booktabs,enumitem,pifont,ulem,fancybox,diagbox,cancel}
\usepackage{pstricks,pstricks-add,pst-plot,pst-tree,pst-node,pst-text,pst-func,pst-math,pst-bspline,pst-3dplot,pst-eucl,pst-solides3d}
\let\pscalculate\undefined
\usepackage{pst-bezier}
\usepackage{tikz,pgfplots,tkz-tab}
\usetikzlibrary{arrows,arrows.meta,patterns,calc,positioning,trees,decorations.markings,shapes.geometric,shapes.misc,shapes.arrows,intersections,decorations.pathmorphing,decorations.pathreplacing,backgrounds,shadows,babel}
\pgfplotsset{compat=1.18}
\usepgfplotslibrary{fillbetween}
\usepackage[np]{numprint}
\usepackage{listings,makecell,pas-tableur,twoopt}
\let\barre\undefined
\usepackage[pstricks]{tablvar}
\definecolor{darkgreen}{RGB}{0,100,0}
\providecommand{\red}{\color{red}}
\providecommand{\blue}{\color{blue}}
\providecommand{\green}{\color{green}}
\providecommand{\black}{\color{black}}
\providecommand{\white}{\color{white}}
\providecommand{\gray}{\color{gray}}
\providecommand{\cyan}{\color{cyan}}
\providecommand{\magenta}{\color{magenta}}
\providecommand{\yellow}{\color{yellow}}
\providecommand{\violet}{\color{violet}}
\providecommand{\orange}{\color{orange}}
\providecommand{\euro}{\texteuro}
\providecommand{\EUR}{\texteuro}
\providecommand{\ieme}{\textsuperscript{e}}
\providecommand{\degres}{\ensuremath{^\circ}}
\providecommand{\degree}{\ensuremath{^\circ}}
\providecommand{\vv}[1]{\overrightarrow{#1}}
\providecommand{\Coord}[1]{#1}
\providecommand{\decosix}{}
\providecommand{\starredbullet}{\ensuremath{\bullet}}
\pagestyle{empty}
\setlength{\parindent}{0pt}
\usepackage{adjustbox}\DeclareUnicodeCharacter{2236}{:}\DeclareUnicodeCharacter{2190}{\ensuremath{\leftarrow}}
\providecommand{\R}{\mathbb{R}}
\providecommand{\N}{\mathbb{N}}
\providecommand{\Z}{\mathbb{Z}}
\providecommand{\Q}{\mathbb{Q}}
\providecommand{\D}{\mathbb{D}}
\providecommand{\C}{\mathbb{C}}
\providecommand{\e}{{\mathrm{e}}}
\providecommand{\vect}[1]{\overrightarrow{#1}}
\providecommand{\vectt}[1]{\overrightarrow{\mathrm{#1}}}
\providecommand{\vvt}[1]{\overrightarrow{\mathrm{#1}}}
\providecommand{\np}[2][]{{#2~#1}}
\providecommand{\vv}[1]{{\overrightarrow{#1}}}
\renewcommand{\d}{{\mathrm{d}}}
\providecommand{\degres}{{^\circ}}
\providecommand{\degre}{{^\circ}}
\providecommand{\blue}{\color{blue}}
\providecommand{\red}{\color{red}}
\providecommand{\enskip}{\quad}
\providecommand{\llbracket}{\lbrack\!\lbrack}
\providecommand{\rrbracket}{\rbrack\!\rbrack}
\providecommand{\euro}{{\text{€}}}
\providecommand{\no}{{\text{n°}}}
\providecommand{\interval}[2]{{[#1;#2]}}
\providecommand{\intervFO}[2]{{[#1;#2[}}
\providecommand{\dotfill}{\ldots}

\providecommand{\ds}{\displaystyle}
\providecommand{\textding}[1]{#1}
\providecommand{\Oij}{\(\mathrm{O};\vec{\imath},\vec{\jmath}\)}
\providecommand{\Oijk}{\(\mathrm{O};\vec{\imath},\vec{\jmath},\vec{k}\)}
\providecommand{\Ouv}{\(\mathrm{O};\vec{u},\vec{v}\)}
\providecommand{\intervFF}[2]{\ensuremath{[#1;#2]}}
\renewcommand{\euro}{\ifmmode\text{\texteuro}\else\texteuro\fi}
\DeclareUnicodeCharacter{00B0}{\ifmmode^{\circ}\else\textdegree\fi}
\DeclareUnicodeCharacter{20AC}{\ifmmode\text{\texteuro}\else\texteuro\fi}
% Énoncés : Éducation nationale. Transcription et corrigés : APMEP, Denis Vergès et contributeurs. Découpage et classement : CoopMaths, Sébastien Lozano, Éric Elter et contributeurs. Adaptation de la présentation : Maths Explorer.
% CC BY-SA : https://coopmaths.fr/www/about/
% https://coopmaths.fr/alea/static/bac/2021/tex/bac_2021_06_metropole_5_cor.tex
\begin{document}
\noindent\textit{APMEP / CoopMaths. CC BY-SA. Présentation adaptée par Maths Explorer.}\par\medskip
\vspace{0.5cm}










On considère l’équation différentielle  $(E)\quad  y'=  y+ 2 x \e^x$
 

\smallskip

\begin{enumerate}
\item 
De $u(x) = x^2\e^x$, on déduit que $u'(x) = 2x\e^x +  x^2\e^x =  \e^x\left(x^2 + 2x \right) = x(x + 2)\e^x$.

Donc $u$ solution de $(E)$ si et seulement si :

$u' = u + 2x\e^x \iff 2x\e^x +  x^2\e^x = x^2\e^x + 2 x \e^x$ qui est vraie : $u$ est 
une solution particulière de $(E)$.
\item Soit $g(x)=f(x) - u(x)$



\smallskip

\begin{enumerate}
\item  
$f$ est solution de l'équation différentielle $(E)$ si et seulement si :

$f'(x) = f(x) + 2x\e^x \quad (1)$.

Or $g(x) = f(x) - u(x) \iff f(x) = g(x) + u(x)$, d'où on déduit, les deux fonctions étant dérivables sur $\R$ : $f'(x) = g'(x) + u'(x)$.

L'égalité (1) devient : $g'(x) + u'(x) = g(x) + u(x) + 2x\e^x \quad (2)$.

Or on a vu dans la question précédente que $u'(x) = u(x) + 2x\e^x$

L'équation (2) devient donc : $g'(x) = g(x)$, ce qui signifie que la fonction $g$ est solution de l'équation différentielle : $y' = y$.



\item 
On sait que les solutions de l'équation différentielle $y' = y$ sont les fonctions définies par $x \longmapsto K\e^x$, \, $K \in \R$.

Donc on a $g(x) = K \e^x, \, K \in \R$ et $f(x) = K \e^x + 2x\e^x$.

Les solutions de l'équation $(E)$ : $f(x) = (K + 2)\e^x$, \, $K \in \R$.
\end{enumerate}
\item Étude de la fonction $u $
	\begin{enumerate}
		\item 
On a $u'(x) = x(x + 2)\e^x$. Comme $\e^x > 0$, quel que soit $x \in \R$, le signe de $u'(x)$ est celui du trinôme $x(x + 2)$ qui a pour racines $- 2$ et $0$.

On sait que ce trinôme est positif, sauf entre les racines :

$u'(x) > 0$ sur $]- \infty~;~-2[ \cup [0~;~+ \infty[$ ;

$u'(x) < 0$ sur $]-2~;~0[$ ;

$u'(- 2) = u'(0) = 0$.
		\item 
		
De la question précédente il suit que $u$ est croissante sauf sur $]-2~;~0[$ où elle est décroissante, $u(- 2) = 4\e^{-2}$ et $u(0)  = 0$ étant les deux extremums de la fonction sur $\R$.

\begin{center}
\psset{unit=1cm}
\begin{pspicture}(10,2.75)
\psframe(10,2.75)\psline(0,2)(10,2)\psline(1,0)(1,2.75)
\uput[u](0.5,2){$x$}\uput[u](1.5,2){$- \infty$}\uput[u](4,2){$- 2$}\uput[u](7,2){$0$}\uput[u](9.5,2){$+\infty$}
\uput[d](4,2){$4\e^{-2}$}\uput[u](7,0){$0$}
\rput(0.5,1){$u$}
\psline{->}(1.5,0.5)(3.5,1.5)\psline{->}(4.5,1.5)(6.5,0.5)\psline{->}(7.5,0.5)(9.5,1.5)
\end{pspicture}
\end{center}
		\item 
$u'$ est un produit de fonctions dérivables sur $\R$, donc elle est dérivable sur $\R$: 

$u''(x) = (2x + 2)\e^x + \left(x^2 + 2x\right)\e^x = \e^x\left(x^2 + 4x + 2\right)$.

Comme $\e^x > 0$, quel que soit $x \in \R$, le signe de $u''(x)$ est celui du trinôme $x^2 +4x + 2 = (x + 2)^2 - 4 + 2 = (x + 2)^2 - 2 =  (x + 2)^2 - \left(\sqrt{2}\right)^2 = \left(x + 2 + \sqrt{2} \right)\left(x + 2 - \sqrt{2} \right)$.

Les racines de ce trinôme sont donc $-\sqrt{2} - 2$ et  $-\sqrt{2} + 2$.

Le trinôme donc $u''(x)$ sont négatifs entre les racines.

Conclusion : la fonction est concave sur l'intervalle $\left]- \sqrt{2} - 2~;~+ \sqrt{2} - 2\right[$.

\begin{center}
\psset{unit=1cm}
\begin{pspicture*}(-6,-0.6)(2,4)
\psaxes[linewidth=1.25pt,labelFontSize=\scriptstyle]{->}(0,0)(-6,0)(2,4)
\psplot[plotpoints=2000,linewidth=1.25pt,linecolor=red]{-6}{2}{x dup mul 2.71828 x exp mul}
\uput[d](-3.414,0){\scriptsize \blue $- \sqrt{2} - 2$}
\uput[d](-0.586,0){\scriptsize \blue $\sqrt{2} - 2$}
\psline[linewidth=1.75pt,linecolor=blue]{]-[}(-3.414,0)(-0.586,0)
\end{pspicture*}
\end{center}
	\end{enumerate}
\end{enumerate}

\end{document}
